% Generado por python -m analysis.capitulo5 (no editar a mano).
\begin{table}[htbp!]
    \centering
    \caption{Dos caídas al azar contra dos caídas dirigidas al nodo más crítico}
    \label{tab:cap5-robustez}
    \footnotesize
    \setlength{\tabcolsep}{3pt}
    \begin{tabular}{|>{\raggedright\arraybackslash\hspace{0pt}}p{2.8cm}|>{\raggedright\arraybackslash\hspace{0pt}}p{2.1cm}|>{\raggedright\arraybackslash\hspace{0pt}}p{2.2cm}|>{\raggedright\arraybackslash\hspace{0pt}}p{2.2cm}|>{\raggedright\arraybackslash\hspace{0pt}}p{2.2cm}|>{\raggedright\arraybackslash\hspace{0pt}}p{2.0cm}|}
        \hline
        \textbf{Experimento} & \textbf{Malla partida (s)} & \textbf{Entrega de datos} & \textbf{Entrega con camino} & \textbf{Cobertura media} & \textbf{Rutas fantasma} \\
        \hline
        Fallos al azar & 42 $\pm$ 57 & 0.502 $\pm$ 0.114 & 0.686 $\pm$ 0.069 & 0.925 $\pm$ 0.067 & 0.020 $\pm$ 0.013 \\
        \hline
        Fallos dirigidos & 180 & 0.284 $\pm$ 0.088 & 0.696 $\pm$ 0.087 & 0.917 & 0.052 $\pm$ 0.003 \\
        \hline
    \end{tabular}

    \vspace{0.4em}\parbox{14.7cm}{\scriptsize Media $\pm$ desviación estándar de 10 corridas de 240 s con 6 flujos de datos, caídas a los 60 y 120 s (lote \texttt{lotes/robustez.json}).}
\end{table}
